% sections/introduction.tex — Introduction
\section{Introduction}
\label{sec:introduction}

\subsection{The Gap Between Actors and GPUs}

The actor model, introduced by Hewitt et~al.~\citep{hewitt1973actors}
and formalized by Agha~\citep{agha1986actors}, takes a position on
concurrency that has proved durable: an actor is a unit of
computation that has private state, processes messages one at a
time from a logically unbounded mailbox, and may create new actors
or send messages to known actors as a result of processing.  No
shared memory, no locks, no threads as a primitive.  The model is
the foundation of Erlang/OTP~\citep{armstrong2007erlang},
Akka~\citep{akka2010}, Pony~\citep{clebsch2015pony}, and—at a
larger scale—the design of fault-tolerant
microservices~\citep{vernon2015reactive}.

GPU programming has evolved along an orthogonal axis.  CUDA and its
successors~\citep{nickolls2008cuda} expose the device as a
coprocessor: the host launches a \emph{kernel} (a function with
implicit parallelism), the kernel runs to completion across a
two-dimensional grid of thread blocks, and any state that needs to
survive must be explicitly copied back to host memory.  The unit of
abstraction is the kernel launch, not a long-lived computation; the
unit of state is the buffer, not the actor.

These two threads of work have been almost entirely
disjoint.  Cross-cutting work exists—\emph{persistent
kernels}~\citep{gupta2012persistent,aila2009raytrace} avoid relaunch
overhead by keeping a single kernel running and feeding it work via
device-side queues; producer–consumer pipelines~\citep{steinberger2014whippletree}
build dataflow graphs on the device; heterogeneous actor frameworks
such as OpenCL Actors~\citep{hiesgen2018opencl} integrate GPU
kernels into host-resident actor runtimes; and more recent
device-side task runtimes (Mustard~\citep{mustard2025}) and
task-based programming models for Hopper-class GPUs
(Cypress~\citep{cypress2025}) attack host-orchestration overhead
from the task-graph side.  In each of those directions the
abstraction is a scheduling or dispatch mechanism rather than a
formal account of a long-lived, device-resident actor with its own
lifecycle, restart policy, causal clock, tenant tag, and
supervisor.  To our knowledge, no existing work combines
device-resident lifecycle, supervision, causal clocks, tenant
routing, and cross-GPU migration under a single formal account.

\subsection{What Hopper Made Possible}

Three hardware features in NVIDIA Hopper (compute capability 9.0)
remove the technical obstacles that previously made persistent
GPU actors awkward to formalize, let alone implement:

\begin{enumerate}
\item \textbf{Thread Block Clusters and Distributed Shared
Memory~(DSMEM)}~\citep{nvidia2022hopper} extend the GPU memory
hierarchy with a new tier between shared memory (per-block) and
global memory (per-device).  Blocks within a cluster can address
each other's shared memory at \(\sim\)30 cycles of latency,
roughly an order of magnitude better than a global-memory
round-trip.  This is the substrate that makes intra-cluster
kernel-to-kernel (K2K) messaging viable as a first-class operation
rather than a host-bounced affair.

\item \textbf{Cooperative kernel launches} with
\texttt{cuLaunchCooperativeKernel} and \texttt{cluster.sync()}
permit grid-wide synchronization within a single kernel invocation.
A persistent kernel can therefore run a producer–consumer loop with
a synchronization barrier at the end of each iteration, without
returning to the host.  This makes the actor's tick loop a real
computational primitive rather than a host-orchestrated
mini-launch sequence.

\item \textbf{Stream-ordered asynchronous
allocation}~(\texttt{cuMemAllocAsync}, \texttt{cuMemFreeAsync})
permits the actor to allocate and free memory inside its own
execution stream without the global-synchronization penalty of
classical \texttt{cuMemAlloc}.  We measured a 116.9\(\times\)
allocation latency improvement over the classical path on H100
(878\,ns versus 102.6\,\(\mu\)s for 256-byte allocations); the
formal model treats memory allocation as an asynchronous,
stream-ordered side effect, exactly mirroring this primitive.
\end{enumerate}

In combination, these features make the design practical enough to
formalize and evaluate directly: Hopper is, to our knowledge, the
first NVIDIA architecture on which a long-lived actor with a
formally-stated lifecycle, supervisor-driven restart, and causal
time can plausibly live on the device without the host orchestrating
every transition.  We are not claiming that Hopper is the first
\emph{conceivable} substrate for such a model; we note only that
sub-microsecond K2K, sub-microsecond cluster synchronization, and
stream-ordered allocation bring device-resident actors within the
cost envelope that actor runtimes typically require of their
substrates.

\subsection{Contribution}

We claim seven contributions, presented in the order they unfold in
the paper.  Each is labeled with what it actually is: a
\emph{definition} (formal statement of a calculus or semantics), an
\emph{analytical proof or proof sketch} (pencil-and-paper argument
over the calculus), a \emph{bounded mechanization} (the corresponding
\TLA{} module model-checked by TLC under specified finite parameter
ranges), or an \emph{empirical measurement} (benchmark numbers with
reported confidence intervals).  These modes are complementary and
we are careful not to conflate them:

\begin{enumerate}
\item \textbf{We define} a small-step operational semantics for the
persistent GPU actor lifecycle (\Cref{sec:model}).  The semantics is
a six-state transition system
(\Dormant{}\,\(\to\)\,\Initializing{}\,\(\to\)\,\Active{}\,\(\to\)\,
\Draining{}\,\(\to\)\,\Terminated{}; \Failed{} as a fault attractor)
defined over an explicit memory model that respects the GPU's SIMT
execution semantics.  \textbf{We prove} lifecycle soundness
(\Cref{thm:lifecycle-soundness}) analytically, and \textbf{we
mechanize} the semantics in \TLA{} and discharge the same invariant
under bounded TLC model checking.

\item \textbf{We define} a K2K messaging calculus (\Cref{sec:k2k})
covering message origination, routing, delivery, and dead-letter
disposition for both intra-block and cross-block (DSMEM) channels.
The calculus is parameterized by a backpressure policy
(\(\mathsf{RejectAndDLQ}\) versus \(\mathsf{DropOldest}\)) and an
optional idempotency cache.  \textbf{We prove} (analytically, plus
bounded mechanization) no-loss accounting for accepted envelopes
(\Cref{thm:no-message-loss}), per-(source,\,destination) FIFO
delivery (\Cref{thm:fifo-per-sender}), and—conditional on an
explicitly-stated cache-sizing hypothesis—at-most-once observable
delivery (\Cref{thm:no-duplicate-delivery}).

\item \textbf{We define} a Hybrid Logical Clock integration that
gives causal ordering on a single device (\Cref{sec:supervision}).
\textbf{We prove} (\Cref{thm:hlc-restart-monotone}) that
supervisor-driven restart preserves HLC monotonicity: if an actor
fails at HLC time \(t\) and is restarted from a snapshot taken at
\(t' < t\), the new instance issues no message with HLC time
\(\le t\).  Restart may roll back \emph{state}, but not
\emph{time}.  To our knowledge this is the first formal statement
of HLC monotonicity across an actor restart; the proof is
analytical, and the per-node monotonicity and event-count
invariants are additionally bounded-mechanized in
\texttt{hlc.tla}.

\item \textbf{We define} a tenant-isolation model (\Cref{sec:tenancy})
in which K2K sub-brokers are keyed on a two-tuple audit boundary
\((\mathit{org\_id}, \mathit{engagement\_id})\).  \textbf{We prove}
(\Cref{thm:cross-tenant-isolation}) that the supervisor's routing
function admits no execution trace in which a message originating
in tenant~\(T_1\) is delivered to an actor registered in
tenant~\(T_2 \ne T_1\).  This is a routing-level structural
guarantee under the model—not an end-to-end confidentiality claim
against an adversary who can subvert the runtime or hardware—but it
is at the right granularity for audit (\Cref{sec:tenancy:limits});
\textbf{we mechanize} the same property under bounded TLC.

\item \textbf{We define} a multi-GPU migration protocol
(\Cref{sec:migration}) that relocates a live actor from source
GPU~\(g_s\) to target GPU~\(g_t\) through a five-phase machine
(\state{Running} \(\to\) \state{Quiesced} \(\to\)
\state{Transferring} \(\to\) \state{Swapping} \(\to\)
\state{Complete}).
\textbf{We sketch} the safety conjunction of atomic-swap,
no-message-loss, FIFO-across-cutover, and state-consistency
(\Cref{thm:migration-safety}), and state the path-selection
property that NVLink-paired actors never fall back to the host
stage (\Cref{thm:fast-path-selection}).  The full analytical
treatment is deferred to a companion paper;
\textbf{we mechanize} both properties in \TLA{} and discharge them
under bounded TLC.  The HLC monotonicity argument extends to
migration via the \(\hlc.\mathsf{merge}\) operation at restore.

\item \textbf{We mechanize} the entire formal apparatus as six
\TLA{} modules under
\texttt{docs/\allowbreak verification/} and model-check them with
TLC under bounded configurations (\Cref{sec:eval:tlc},
\Cref{tab:tlc-stats}); all six pass with no counterexamples within
their stated bounds.  Running the pass caught a real
specification-level bug in the migration checksum invariant
(\Cref{sec:disc:tlc-bug}) before it reached the paper—a concrete
instance of the verification apparatus doing engineering work rather
than ceremonial proof dressing.

\item \textbf{We measure} the runtime on NVIDIA H100~NVL silicon
(\Cref{sec:evaluation}).  A 1{,}000-command burst submitted
through mapped memory costs 0.18\,\(\mu\)s of host wall-clock
\emph{amortized over the burst through the CPU's write-combining
buffer}, against 1{,}583.1\,\(\mu\)s for the same burst issued as
1{,}000 \texttt{cuLaunchKernel} calls---8{,}698\(\times\) less
host time submitting the burst on this path, under the control-path
/ work-submission comparison described in \Cref{sec:eval:injection}
(not a runtime-vs-runtime execution-speed claim).  K2K messaging yields
72--83\,ns round-trip latency at payloads up to 1~KiB through the
intra-block lock-free queue, with 5.10\,M~messages/second
sustained throughput at coefficient of variation 0.66\% over four
60-second trials.  Cluster-level synchronization runs at
0.628\,\(\mu\)s (2.98\(\times\) faster than grid-wide
synchronization); stream-ordered asynchronous allocation is
116.9\(\times\) faster than classical \texttt{cuMemAlloc}; and
NVLink P2P reaches 257.78\,GB/s unidirectional at 16~MiB on a
2\(\times\) H100~NVL system, which is 85.9\% of the 12-link
NVLink4 bundle's 300\,GB/s unidirectional peak (equivalently,
43.0\% of its 600\,GB/s bidirectional peak).
We also report the 4.2\(\times\) slowdown of a block-actor
WaveSim3D stencil relative to a hand-tuned kernel
(\Cref{tab:apps}): this is the price of actor semantics over a
fused baseline, which we take to be among the more direct
indicators of the cost of the abstraction.  All measurements were collected under locked
GPU clocks, exclusive compute mode, and a statistical protocol
calibrated against the reproducibility criteria of
\citet{georges2007rigorous} and~\citet{hoefler2015scientific}.
\end{enumerate}

\subsection{Audience and Scope}

The paper is written for two audiences.  For the \emph{programming
languages} reader, \Cref{sec:model,sec:k2k,sec:supervision,sec:tenancy}
develop the formal model in the style of small-step operational
semantics with proofs at the level of detail customary in the
field.  For the \emph{GPU systems} reader, \Cref{sec:implementation}
maps the formal model onto the \system{} runtime and discusses the
engineering choices (CUDA driver-API surface, \texttt{cuMemAllocAsync}
integration, NVTX profiling instrumentation) needed to realize each
semantic primitive.  \Cref{sec:evaluation} presents the
H100~NVL benchmarks and validates the numbers stated above.

What the paper does \emph{not} cover: a complete account of the
GPU memory-consistency model (we adopt the axiomatic model of
\citet{lustig2014abstraction} and constrain our proofs to the
operations that model defines), nor a comprehensive performance
study against alternative GPU programming
abstractions~(CUDA Graphs, Triton, JAX); we focus on the
formal model and its mechanically-checkable properties, with
benchmarks sized to validate the model rather than to claim
superiority across all use cases.
